\subsection{Dataset Description}
\label{subsec:data_description}

\begin{sloppypar}

Three publicly available Kaggle datasets are used, all framed around the same
binary target: whether a respondent or patient has (or has had) heart disease
or a heart attack. The datasets differ substantially in size, feature type, and
data-quality characteristics, which is intentional: comparing AdaBoost, XGBoost,
and LightGBM across three heterogeneous datasets, rather than a single one,
lets RQ1--RQ3 be answered under different sample sizes, feature compositions,
and class-imbalance regimes rather than for one dataset in isolation. Each
dataset is loaded directly from its raw CSV file under \texttt{data/raw/} via a
single shared \texttt{load\_raw\_dataset()} function, with no dataset-specific
loading logic.

\begin{table}[H]
\centering
\begin{minipage}{\linewidth}
\centering
	\refstepcounter{table}
	\begin{captionbelowboxaivn}
		\captable{Raw shape, feature-type composition, missingness, and target positive rate for the three datasets before any preprocessing.}
		\label{tab:dataset_overview}
	\end{captionbelowboxaivn}

	\renewcommand{\arraystretch}{1.1}
	\small

	\begin{tabularx}{\linewidth}{@{}
	>{\raggedright\arraybackslash}X
	>{\centering\arraybackslash}X
	>{\centering\arraybackslash}X
	>{\centering\arraybackslash}X
	>{\centering\arraybackslash}X
	>{\centering\arraybackslash}X@{}}
\textbf{Dataset} & \textbf{Rows} & \textbf{Numeric} & \textbf{Categorical} & \textbf{Missing} & \textbf{Positive rate} \\
\hline
Dataset 1 & 445{,}132 & 6 & 34 & 38/40 columns & 5.64\% \\
Dataset 2 & 918 & 6 & 5 & none & 55.34\% \\
Dataset 3 & 253{,}680 & 7 & 14 & none & 9.42\% \\
		\end{tabularx}
\end{minipage}
\end{table}

\textbf{Dataset 1 --- Personal Key Indicators of Heart Disease (2022):} sourced
from the Kaggle 2022 BRFSS release \cite{kaggle2022heartindicators} and
distributed as \texttt{heart\_2022\_with\_nans.\allowbreak csv}. This variant is used
deliberately, rather than the pre-cleaned no-NaN release, so that the
preprocessing pipeline demonstrates a realistic missing-value-handling
workflow. The raw table has 445{,}132 rows and 40 columns: 34 categorical
columns (\texttt{State}, \texttt{Sex}, \texttt{GeneralHealth}, and most of the
Yes/No health-history flags such as \texttt{HadHeartAttack}, \texttt{HadDiabetes},
\texttt{HadStroke}) and 6 numeric columns (\texttt{PhysicalHealthDays},
\texttt{MentalHealthDays}, \texttt{SleepHours}, \texttt{HeightInMeters},
\texttt{WeightInKilograms}, \texttt{BMI}). The target column,
\texttt{HadHeartAttack}, is a Yes/No self-report that is mapped to a binary
0/1 label during preprocessing; before cleaning, 93.67\% of respondents
(416{,}959) answered "No", 5.64\% (25{,}108) answered "Yes", and 0.69\% (3{,}065)
left the field blank, giving a roughly 17:1 imbalance between the majority and
minority classes once the missing-target rows are dropped.

Thirty-eight of the 40 columns contain at least one missing value; only
\texttt{State} and \texttt{Sex} are fully populated. Table
\ref{tab:d1_missing} lists the six most-affected columns, all
vaccination/screening-history fields that respondents frequently skip or do
not know the answer to. The core numeric health metrics have moderate
missingness (\texttt{BMI} 10.96\%, \texttt{WeightInKilograms} 9.45\%,
\texttt{HeightInMeters} 6.44\%), while the disease-history Yes/No flags and
core demographics are the most complete, at well under 1\% missing --- a
pattern typical of self-reported survey data, where harder-to-recall or less
commonly asked questions accumulate more gaps.

\begin{table}[H]
\centering
\begin{minipage}{0.8\linewidth}
\centering
	\refstepcounter{table}
	\begin{captionbelowboxaivn}
		\captable{The six most-missing columns in Dataset 1 (raw data, before imputation).}
		\label{tab:d1_missing}
	\end{captionbelowboxaivn}

	\renewcommand{\arraystretch}{1.3}
	\footnotesize

	\begin{tabularx}{\linewidth}{@{}
        >{\raggedright\arraybackslash}X
        >{\centering\arraybackslash}X@{}}
\textbf{Column} & \textbf{Missing rate} \\
\hline
TetanusLast10Tdap & 18.54\% \\
PneumoVaxEver & 17.31\% \\
HIVTesting & 14.86\% \\
ChestScan & 12.59\% \\
CovidPos & 11.40\% \\
HighRiskLastYear & 11.37\% \\
	\end{tabularx}
\end{minipage}
\end{table}

\textbf{Dataset 2 --- Heart Failure Prediction:} sourced from the fedesoriano
Kaggle release \cite{kaggle2020heartfailure} and distributed as
\texttt{heart.csv}. This is the smallest and already mostly-clean table used
in this study: 918 rows and 12 columns, combining 5 clinical/demographic
categorical columns (\texttt{Sex}, \texttt{ChestPainType}, \texttt{RestingECG},
\texttt{ExerciseAngina}, \texttt{ST\_Slope}) with 6 numeric columns
(\texttt{Age}, \texttt{RestingBP}, \texttt{Cholesterol}, \texttt{FastingBS},
\texttt{MaxHR}, \texttt{Oldpeak}) and the pre-encoded binary target
\texttt{HeartDisease}. There are zero missing values across all 918 rows,
so the dataset appears complete on the surface, but two of its numeric
columns use \texttt{0} as an undocumented placeholder for "not measured"
rather than a genuine reading: exactly 1 row has \texttt{RestingBP == 0} and
172 rows (18.7\%) have \texttt{Cholesterol == 0}, far too many to be real
zero-cholesterol readings given that a cholesterol or blood-pressure value of
0 is not physiologically possible. \texttt{DatasetConfig.\allowbreak zero\_as\_missing\_columns}
records both columns so that the preprocessing pipeline recodes these
placeholder zeros as \texttt{NaN} before imputation rather than taking them at
face value. Unlike Dataset 1, the target here is close to balanced: 508
patients (55.34\%) have heart disease and 410 (44.66\%) do not, which is
typical of a clinical dataset curated specifically for a heart-disease
classification task rather than drawn from a general population survey.

\textbf{Dataset 3 --- Heart Disease Health Indicators (BRFSS 2015):} sourced
from the alexteboul Kaggle release \cite{kaggle2015brfssindicators} and
distributed as \texttt{heart\_disease\_health\_indicators\_BRFSS2015.\allowbreak csv}. The
raw table has 253{,}680 rows and 22 columns, all already numerically encoded
as \texttt{float64} -- including conceptually binary flags (\texttt{HighBP},
\texttt{Smoker}, \texttt{Sex}) and small ordinal scales
(\texttt{GenHlth} 1--5, \texttt{Age} 1--13 brackets, \texttt{Education} 1--6,
\texttt{Income} 1--8). Seven columns (\texttt{BMI}, \texttt{GenHlth},
\texttt{MentHlth}, \texttt{PhysHlth}, \texttt{Age}, \texttt{Education},
\texttt{Income}) are treated as numeric and the remaining 14 as categorical.
There are zero missing values anywhere in the dataset, making it the cleanest
of the three in terms of completeness. However, 23{,}899 rows (9.42\%) are
exact duplicates of another row across all 22 columns -- plausible here
because every feature is binary or a small discrete scale, so two distinct
respondents can easily share an identical combination of answers. This is a
data-quality issue distinct from missingness, and \texttt{DatasetConfig}
flags Dataset 3 with \texttt{drop\_\allowbreak duplicate\_\allowbreak rows=True} so these rows are
removed before the target is cleaned (Section~\ref{subsec:data_preprocessing}).
The target, \texttt{HeartDiseaseorAttack}, is heavily imbalanced on the raw
data: 229{,}787 respondents (90.58\%) report no heart disease/attack history
and 23{,}893 (9.42\%) do, a roughly 1:9.6 minority-to-majority ratio, similar
in shape to Dataset 1 but somewhat less extreme.

Across all three datasets, the target definition, raw imbalance, and
data-quality quirks summarized above directly motivate the dataset-specific
preprocessing decisions described in Section~\ref{subsec:data_preprocessing}:
the class-imbalance policy (F1), the zero-as-missing recoding restricted to
Dataset 2, and the duplicate-row removal restricted to Dataset 3.
\end{sloppypar}
